\chapter{Conclusion}
\label{ch:conclusion}

This report set out to build and evaluate a persona-drift-gated retrieval architecture over a typed-memory persona representation. The architecture as originally designed (v0.3 of the project's PRD) used a persona-vector projection as both a rerank scorer (M2) and a drift-gate signal (M3), with the persona vector extracted by the Anthropic persona-vectors methodology and validated on contrast prompts. The architecture as finally evaluated (v3.1) drops M2, reverts M3 to an LLM-as-judge gate, and retains typed retrieval with per-turn identity grounding as the affirmative architectural contribution (M1).

The reshaping happened because Experiment 2 refuted the architectural assumption that the contrast-trained persona-vector direction would transfer to inference-time drift detection in multi-turn dialogue at the scale and quantisation the project tested. Experiment 1 replicates the persona-vectors extraction cleanly on Gemma-2-9B-Instruct at 4-bit NF4 with the Dubanowska defensive controls intact (per-layer test \auroc{} 1.000 on all three personas at all four swept layers; random-feature control 0.58 to 0.65, below the 0.70 weak floor; shuffled-label at chance). Experiment 2 then tests transfer across the 24 generation-scope cells on Gemma and Llama plus the 12 prompt-scope cells on Gemma. The maximum absolute projection delta between in-persona and drifting conversations is 0.103, in the wrong direction, against a pre-registered $+0.30$ proceed threshold. The signal that M3's gate and M2's reranker would have keyed on is missing by an order of magnitude across every cell.

What survives the refutation is a typed-memory mechanism (M1) and a drift-gated mechanism with an LLM-judge gate (M3 v3.1). Evaluated against a vanilla-RAG floor (B1) and a RoleGPT-level prompt baseline (B2) on a custom counterfactual-retrieval probe suite of 90 multi-turn conversations across three role personas and three probe types, scored by MiniCheck, an adapted SYCON, and a three-judge PoLL panel, the four pipelines show two consistent patterns. B1 collapses (PoLL persona-adherence 2.54 against the cluster's 4.5 to 4.7), demonstrating that persona must be encoded somewhere in the pipeline at all. B2, M1, and M3 cluster within 0.20 PoLL points of each other on every probe type, including the Type B counterfactual-injection probes the mechanisms were designed to address. The targeted-hypothesis prediction that retrieval-side conditioning would beat prompt-only persona on Type B is not supported. B2 leads on Type B.

The mechanism cluster's robustness is the contribution of the diagnostic ablations. An oracle drift gate that fires by construction on the known drift turns does not unlock headroom (PoLL persona-adherence 4.57, inside the cluster). M1 without per-turn identity grounding produces the same headline score (4.57). The 4-bit-to-fp16 precision sweep produces a $\Delta$-of-$\Delta$ of $-0.142$, inside the pre-registered $\pm 0.15$ envelope. Three confounds isolated, three confirmations that the cluster sits where it sits because the binding constraint at the 9B parameter scale is model capability, not retrieval architecture, gate calibration, or quantisation level.

The methodological observation that comes out of the two-experiment design is that when a published validation of an inference-time signal uses one regime (contrast prompts, single-turn) and the proposed architectural application uses another regime (multi-turn dialogue), the transfer is a real empirical question and should be tested before commitment. The cost is small. The architectural saving, if the transfer fails, is large. The 36 forward passes per (backbone, scope) cell that Experiment 2 ran are enough to surface a robust refutation. The same kind of sweep would have caught the assumption in the v0.3 design before M2 and M3 were scoped around it.

The contributions are three. A replication of persona-vector extraction at a (backbone, quantisation) cell the published literature did not cover (Experiment 1). A robust negative result on persona-vector transfer to multi-turn drift detection on two production-scale backbones in two extraction regimes (Experiment 2). An evaluation of what survives the refutation, on a custom counterfactual-retrieval probe suite, with three diagnostic ablations that isolate the cluster's source as model capability rather than retrieval design. The project ships a typed-memory architecture, a persona-vector extraction and caching pipeline, an evaluation harness with the metric stack and the PoLL panel, and a probe suite calibrated to the pre-registered $30\%$ to $70\%$ B2 failure band. Every reported number traces to a run directory in the project repository.

The negative result is the most actionable finding for subsequent work. The persona-vectors line of research has produced clean replications across several model families and quantisation regimes. The architectural extension to multi-turn drift gating, as proposed by this project and presumably by others working in the same intersection of internal-state-conditioned retrieval and persona-vector probing, does not transfer at this scale by the published methodology. Subsequent work either needs a different signal (the cleanest candidate is a probe trained directly on hidden states sampled from genuinely drifting dialogue, rather than from contrast prompts), a larger model scale at which the geometry may differ, or an architectural commitment that does not rely on the projection as a single-scalar gate. The thesis-bridge follow-ups named in Chapter~\ref{ch:limitations} cover the first and the third of these.

What we set out to build is not what we shipped. What we shipped is a smaller architectural claim about typed memory and ID-RAG that the evaluation places inside a tight cluster with a strong prompt baseline, plus a robust empirical finding about an architectural assumption the prior literature did not test. We think the second is the more useful contribution.
